\chapter{Periodic Systems}
\label{ch:periodic}

The periodic code in \code{pycmf.pbc} is \textbf{experimental}. It
reproduces the previous implementation (version 0.2) and the KMP2 code of
\pyscf{} at the first cycle, but it has not yet been compared with published
numbers. Every class of the sub-package emits an \code{ExperimentalWarning}
when it is instantiated.

\begin{lstlisting}[style=python]
from pycmf.pbc import KOBMP2, KO2BMP2, KOBDF, KUOBMP2
from pycmf.pbc import remove_linear_dependency_
\end{lstlisting}

\section{k-point OBMP2: KOBMP2 and KO2BMP2}
\label{sec:kobmp2}

\subsection{Usage}

\begin{lstlisting}[style=python]
from pyscf.pbc import gto, scf
from pycmf.pbc import KOBMP2, KO2BMP2

a0 = 3.5668
cell = gto.M(atom='C 0 0 0; C %f %f %f' % (a0/4, a0/4, a0/4),
             a=[[0, a0/2, a0/2], [a0/2, 0, a0/2], [a0/2, a0/2, 0]],
             basis='gth-szv', pseudo='gth-pade')
mf = scf.KRHF(cell, cell.make_kpts([2, 2, 2])).density_fit().run(conv_tol=1e-10)

mp = KOBMP2(mf).run()            # e_tot, e_corr per unit cell
print(mp.e_tot, mp.band_gap())   # gap of the bands of the correlated Fock matrix

o2 = KO2BMP2(mf, c_os=1.2).run() # c_os has no default
KOBMP2(mf, kpt_symmetry=False)   # full k-mesh, no symmetry
\end{lstlisting}

With the diamond cell above (\code{examples/experimental/kobmp2.py}):

\begin{lstlisting}[style=output]
E(KRHF)    = -10.93208054   gap = 0.67421
E(KOBMP2)  = -11.02688135   gap = 0.73499
without k-point symmetry: 8 k-points instead of 3, same energy: True
E(KO2BMP2, c_os = 1.2) = -11.00812694   gap = 0.72361
\end{lstlisting}

\subsection{Requirements on the mean field}

\begin{itemize}
  \item The reference must be a closed-shell \code{KRHF} object with integer
        occupations. Kohn--Sham objects are rejected. The number of
        occupied bands may depend on the $k$-point.
  \item The system must be three-dimensional.
  \item The correlation part uses Gaussian density fitting. If the mean field
        uses a \code{GDF} or \code{RSGDF} object, it is reused; otherwise
        (for example FFT-based integrals) a \code{GDF} object is built for
        the correlation part, and the result differs from an all-FFT
        calculation by the density-fitting error.
  \item A mean field that uses $k$-point symmetry (\code{KPoints} objects of
        \pyscf{}) is converted to the full mesh.
\end{itemize}

\subsection{Attributes}

\begin{longtable}{@{}p{0.28\textwidth} p{0.17\textwidth} p{0.47\textwidth}@{}}
\caption{Attributes of \code{KOBMP2}.}\label{tab:kattrs}\\
\toprule
Attribute & Default & Meaning \\
\midrule
\endfirsthead
\toprule
Attribute & Default & Meaning \\
\midrule
\endhead
\bottomrule
\endfoot
\code{conv_tol} & \code{1e-8} & convergence of the energy per cell \\
\code{conv_tol_grad} & $\sqrt{\text{\code{conv_tol}}}$ & convergence of the occupied--virtual block of $\fbar$ \\
\code{max_cycle} & \code{100} & maximum number of cycles \\
\code{diis}, \code{diis_space}, \code{diis_start_cycle} & \code{True, 8, 1} & DIIS on the Fock matrices of the irreducible $k$-points \\
\code{bch_order} & \code{2} & 1: half-weighted first commutator \\
\code{c_os}, \code{c_ss} & \code{1, 1} & amplitude scaling; \code{KO2BMP2} sets \code{c_ss = 0} \\
\code{level_shift} & \code{0} & shift of the energy denominators \\
\code{frozen} & \code{None} & number of lowest bands, or list of band indices, without amplitudes \\
\code{kpt_symmetry} & \code{True} & \code{True}: space group and time reversal; \code{'time_reversal'}; \code{False} \\
\code{with_df} & from \code{mf} & GDF or RSGDF object; \code{mp.density_fit(auxbasis=...)} builds another one \\
\code{max_memory} & from \code{mf} & decides whether three-index integrals stay in memory \\
\code{direct_exchange} & \code{True} & take the HF exchange matrices from the pass over the integrals that builds the amplitudes \\
\code{single_precision} & \code{False} & store the occupied--virtual three-index tensors in single precision \\
\end{longtable}

\code{frozen} and \code{kpt_symmetry} can also be given to the constructor,
\code{KOBMP2(mf, frozen=None, kpt_symmetry=None)}. Frozen orbitals are
excluded from the amplitudes at every $k$-point but are still relaxed by the
correlated Fock operator. The frozen set must not cut a degenerate level; in
diamond the two lowest valence bands are degenerate at X, so \code{frozen=1}
is not valid there (a warning is issued).

\code{direct_exchange} gives identical energies and saves one pass over the
integral file per cycle; it applies to a plain \code{KRHF} with GDF and
without MPI. \code{single_precision} halves the memory of the stored tensors;
the arithmetic stays in double precision, and the energy changes by less than
$10^{-9}$~Hartree for diamond gth-dzvp $3\times3\times3$.

The names of version 0.2 are kept as aliases: \code{niter}
(\code{max_cycle}), \code{thresh} (\code{conv_tol}), \code{shift}
(\code{level_shift}) and \code{ene_tot} (\code{e_tot}).

\subsection{Results and methods}

After \code{kernel()}: \code{converged}, \code{e_tot}, \code{e_corr}
(relative to KRHF), \code{e_ref} (HF energy of the OBMP2 determinant), and
\code{mo_energy}, \code{mo_coeff} and \code{mo_occ} for all $k$-points of the
mesh (lists over $k$-points, occupied orbitals first). \code{fock_hf} holds
the HF Fock matrices in the basis of the final orbitals.

\begin{longtable}{@{}p{0.42\textwidth} p{0.52\textwidth}@{}}
\toprule
Method & Returns \\
\midrule
\endhead
\bottomrule
\endfoot
\code{get_fock()} & correlated Fock matrices (AO basis) at all $k$-points \\
\code{ip_koopmans()}, \code{ea_koopmans()} & minus the valence band maximum and minus the conduction band minimum on the mesh \\
\code{band_gap()} & gap of the eigenvalues of $\fbar$ on the mesh \\
\code{init_from_scf()} & takes the orbitals of the mean field without running OBMP2 (for conventional second-order quasiparticle energies) \\
\code{ipea(nvir=4, kpts_index=None)} & second-order IP/EA corrections; returns the quasiparticle gap \\
\code{self_energy(nvir=2, ...)} & tabulated second-order self-energy of the band edges \\
\code{quasiparticle_gap(scheme, reference)} & (direct gap at the first $k$-point, smallest gap on the mesh) \\
\code{peom_mp2(kpt_v=0, kpts_c=None)} & band edges and gaps from partitioned EOM-MP2 \\
\code{downfold(ncas, nelecas, cas_list=None)} & \code{KOBDF} object \\
\code{density_fit(auxbasis=None, with_df=None)} & sets another density-fitting object \\
\end{longtable}

\subsection{Restart}

To restart from given orbitals (for example a converged calculation with other
scaling factors), set \code{mp.mo_coeff}, \code{mp.mo_energy} and
\code{mp.mo_occ} (lists over the $k$-points of the mesh) before calling
\code{kernel()}.

\section{Algorithm}
\label{sec:kpt-algorithm}

\paragraph{Three-index integrals.} All integrals are products of GDF tensors,
\[
  (ia|jb) = \sum_Q L^Q_{ia}(k_i,k_a)\,L^Q_{jb}(k_j,k_b).
\]
The
occupied--virtual tensors (\emph{rows}, one per $k$-point) are transformed
once per cycle and kept in memory, or in a temporary HDF5 file when they need
more than half of \code{max_memory}. The AO tensors are cached between cycles
as far as the remaining memory allows (up to 80\,\% of \code{max_memory} in
total).

\paragraph{No integral with a general index.} The terms
$\sum \bar t_{iajb}(ia|pb)$ and $\sum \bar t_{iajb}(ia|jp)$ are evaluated
through the intermediate
\begin{equation}
  N^Q_{jb}(k_b) = \sum_{k_i}\sum_{ia} L^Q_{ia}(k_i,k_a)\,\bar t_{iajb},
\end{equation}
accumulated over $k_i$ and contracted once per pair $(k_j,k_b)$ with
$L^Q_{pb}$ and $L^Q_{jp}$, which are built on the fly.

\paragraph{Memory.} Nothing of size $\Nk^2 o^2 v^2$ is stored; the amplitudes
exist for one $k$-point quadruple at a time.

\paragraph{Symmetry.} Every term of the potential is written with its free
indices on $k_j$ (\emph{$j$ terms}) or on $k_b$ (\emph{$b$ terms}), using
$t_{iajb}=t_{jbia}$. $k_j$ then runs over the irreducible $k$-points only:
the $j$ terms are complete, and the $b$ terms are completed by the group
average $v = |G|^{-1}\sum_g g^{-1} P g$ of the partial sum $P$, with the
orbital rotations of \code{pyscf.pbc.symm}. The orbitals are kept at the
irreducible $k$-points and expanded in a symmetry-consistent gauge in every
cycle. If the mean-field solution does not have the symmetry of the mesh, the
full mesh is used. Only symmorphic operations are used (the default of
\pyscf{}).

\paragraph{Self-consistency.} The loop is \code{pycmf.core.scf_loop} with the
irreducible $k$-points as channels: DIIS with the commutator $[F,D]$,
convergence on $\Delta E$ and on the occupied--virtual block of $\fbar$, and
diagonalisation in the space of the current orbitals, so that orbitals
removed because of linear dependencies stay removed.

\section{Resources and performance}
\label{sec:kpt-resources}

\paragraph{Disk.} The three-index file of \pyscf{} holds all pairs of
$k$-points. For diamond gth-dzvp it takes 3.9~GB at $4\times4\times4$ and
14.8~GB at $5\times5\times5$, and \pyscf{} needs about as much again for a
temporary file while it is built. Set \code{mf.with_df._cderi_to_save} to
keep the file and \code{mf.with_df._cderi} to reuse it in a later run.

\paragraph{Memory.} The reading of the AO three-index integrals disappears
when \code{max_memory} is large enough to keep all AO tensors in memory (7.4~GB
for diamond gth-dzvp $4\times4\times4$; the default of 4000~MB holds a third of
them).

\paragraph{Timings.} Table~\ref{tab:ktimings} gives the time per
self-consistent cycle for diamond (two atoms per cell, GTH pseudopotential,
GDF, two threads of a shared virtual machine; timings fluctuate by about
$\pm20$\,\%). The script \code{benchmarks/kpoints_timing.py} reproduces such
measurements on other hardware.

\begin{table}[!htb]
\centering
\small
\caption{Time per KOBMP2 cycle for diamond, including the HF Fock build.}
\label{tab:ktimings}
\begin{tabular}{@{}l c r r r r@{}}
\toprule
Basis, mesh & $k$-points (irred.) & 0.2 & 1.0 full & 1.0 symm. & Memory \\
\midrule
gth-szv, $2^3$  & 8 (3)  & 4.2 s & 0.30 s & 0.17 s & 0.2 GB \\
gth-szv, $3^3$  & 27 (4) & 149 s & 3.7 s  & 2.1 s  & 0.3 GB \\
gth-dzvp, $2^3$ & 8 (3)  & 11 s  & 1.1 s  & 0.8 s  & 0.3 GB \\
gth-szv, $4^3$  & 64 (8) & $\approx$36 min & 33 s & 11.5 s & 0.9 GB \\
gth-dzvp, $3^3$ & 27 (4) & $\approx$7 min  & 27 s & 7.5 s  & 1.7 GB \\
gth-dzvp, $4^3$ & 64 (8) & $\approx$94 min & --   & 88 s   & 3.5 GB \\
\bottomrule
\end{tabular}
\end{table}

The values for version 0.2 have its debugging overhead removed; values marked
$\approx$ are extrapolated with $\Nk^3$. With DIIS, five to seven cycles bring
$|\Delta E|$ below $10^{-8}$~Hartree from HF orbitals; without DIIS seven to
thirteen.

\section{Nearly linearly dependent basis sets}
\label{sec:lindep}

The OBMP2 iterations are less robust than Hartree--Fock when the Bloch basis
is nearly linearly dependent. For diamond with gth-tzvp \pyscf{} removes two
to four functions per $k$-point, but overlap eigenvalues below $10^{-4}$
remain, and with them KOBMP2 does not converge. The function
\code{remove_linear_dependency_(mf, threshold=1e-4)}, called before
\code{mf.kernel()}, makes the mean field work in the space of the overlap
eigenvectors above the threshold:

\begin{lstlisting}[style=python]
from pycmf.pbc import KOBMP2, remove_linear_dependency_
mf = scf.KRHF(cell, kpts).density_fit()
remove_linear_dependency_(mf, 1e-4)
mf.kernel()
mp = KOBMP2(mf).run()
\end{lstlisting}

KOBMP2 keeps that orbital space and converges in five to six cycles, at a
cost of about $5\times10^{-5}$~Hartree in the HF energy. The number of
orbitals then depends on the $k$-point, the mean field is no longer exactly
symmetric, and KOBMP2 falls back to the full mesh.

\section{MPI}
\label{sec:kpt-mpi}

With \code{mpi4py} the loop over $k$-points of the correlation potential is
distributed. The same script runs on every process; rank 0 does the
Hartree--Fock part and broadcasts it, and the partial results are summed with
\code{allreduce}. Output is written by rank 0 only.

\begin{lstlisting}[style=shell]
$ OMP_NUM_THREADS=1 mpiexec -n 4 python kobmp2_mpi.py
\end{lstlisting}

See \code{examples/experimental/kobmp2_mpi.py}. Processes times threads must
not exceed the number of cores. For the small test case BN $3\times3\times3$
gth-szv, two processes reduce three cycles from 10.5~s to 8.2~s with symmetry;
the serial parts (HF build, reading of the integrals) limit the gain for
small cells.

\section{Band gaps and quasiparticle energies}
\label{sec:qp}

The eigenvalues of the correlated Fock matrix are not removal and addition
energies, so \code{band_gap()} is not a quasiparticle gap. Three levels of
theory for the band edges are built on a KOBMP2 calculation.

\subsection{Static correction: IPEA-KOBMP2}

\code{mp.ipea(nvir=4)} adds the second-order corrections
\begin{equation}
  \mathrm{IP}(h,k) = -\varepsilon(h,k) + 2\sum \bar t_{iajh}\,(ia|jh),\qquad
  \mathrm{EA}(l,k) = -\varepsilon(l,k) - 2\sum \bar t_{ialb}\,(ia|lb)
\end{equation}
to the converged eigenvalues (module \code{pbc/kipea.py}), stores
\code{mp.ip} and \code{mp.ea} for all $k$-points and returns the
quasiparticle gap $\min\mathrm{IP}-\max\mathrm{EA}$. The amplitudes carry
the same \code{c_os}, \code{c_ss} scaling as the potential. The corrections
are static and evaluated once; only the band edges are meaningful, because
the denominators can vanish for bands away from the gap.

\subsection{Second-order self-energy}

\code{pbc/kqp.py} tabulates the two parts of the diagonal second-order
self-energy on a frequency grid around each band energy,
\begin{equation}
  A_p(\omega) = \sum_{iab} \frac{W}{\omega + \varepsilon_i - \varepsilon_a - \varepsilon_b}
  \quad(\text{2p1h}),\qquad
  B_p(\omega) = \sum_{iaj} \frac{W}{\omega + \varepsilon_a - \varepsilon_i - \varepsilon_j}
  \quad(\text{2h1p}).
\end{equation}
With HF orbitals the diagonal of the OBMP2 potential is exactly one of them:
$v_{hh}=A_h(\varepsilon_h)$ for occupied and $v_{ll}=B_l(\varepsilon_l)$ for
virtual bands. The eigenvalues of $\fbar$ therefore contain the half of the
static self-energy that opens the gap ($\Sigma_{\mathrm{in}}$), and the IPEA
correction is the other half ($\Sigma_{\mathrm{out}}$). Nothing is counted
twice.

\begin{lstlisting}[style=python]
mp.kernel()
mp.self_energy(nvir=1, nocc=1)          # band edges only
direct, smallest = mp.quasiparticle_gap('dynamic')
\end{lstlisting}

\begin{table}[!htb]
\centering
\small
\caption{Quasiparticle schemes (\code{reference='obmp2'}).}
\begin{tabular}{@{}l l@{}}
\toprule
Scheme & Quasiparticle energy \\
\midrule
\code{'static'}  & $\varepsilon_p + \Sigma_{\mathrm{out}}(\varepsilon_p)$; identical to \code{mp.ipea()} \\
\code{'linear'}  & $\varepsilon_p + Z\,\Sigma_{\mathrm{out}}(\varepsilon_p)$, $Z = 1/(1 - A' - B')$ \\
\code{'dynamic'} & root of $\omega = \varepsilon_p + [\Sigma_{\mathrm{in}}(\omega) - \Sigma_{\mathrm{in}}(\varepsilon_p)] + \Sigma_{\mathrm{out}}(\omega)$ \\
\bottomrule
\end{tabular}
\end{table}

With \code{reference='hf'} together with \code{KOBMP2(mf).init_from_scf()}
the conventional second-order quasiparticle energies on Hartree--Fock are
obtained, $\varepsilon_p + A + B$.

Arguments of \code{self_energy}: \code{nvir} and \code{nocc} are the numbers
of lowest virtual and highest occupied bands that are treated
(\code{nocc=1, nvir=1} is enough for the gap); \code{kpts_index} selects
$k$-points (default: the irreducible ones); \code{grid} gives the frequency
offsets (default 41 points within $\pm0.4$~Hartree or 1.4 times the gap);
\code{scaling} is \code{'linear'} or \code{'hylleraas'}.

\paragraph{Spin-component scaling.} The amplitudes enter the potential
linearly with the factor $c$ and quadratically with $c^2$, so the half inside
the eigenvalues carries $2c-c^2$ per spin component, whereas the IPEA formula
multiplies the other half by $c$ itself. \code{self_energy(scaling='hylleraas')}
evaluates the self-energy with $2c-c^2$, so that both halves are weighted
alike; the default \code{'linear'} reproduces the IPEA numbers. For
$c_{\mathrm{OS}}=c_{\mathrm{SS}}=1$ there is no difference.

The lower-level functions are available directly:
\code{kqp.self_energy(mp, ...)}, \code{kqp.quasiparticle(se, scheme,
reference)}, \code{kqp.band_edges(se, e_qp)} and \code{kqp.gaps(se, e_qp)}.

\subsection{Partitioned EOM-MP2}

\code{pbc/kpeom.py} computes the band edges with partitioned EOM-MP2: the
eigenvalues of $e^{-T_2}\hat He^{T_2}$ in the space of 1h and 2h1p (1p and
2p1h) configurations, with the doubles--doubles block replaced by
orbital-energy differences. Compared with the second-order self-energy, one of
the two integrals of every term is dressed by the amplitudes, which adds the
leading ring (screening), ladder and exchange terms of third order. The method
has no adjustable parameter.

\begin{lstlisting}[style=python]
mp = KOBMP2(mf).run()                 # or KOBMP2(mf).init_from_scf()
res = mp.peom_mp2(kpt_v=0, kpts_c=[0, kc])
res['ip'], res['ea'][kc], res['gap'][kc]     # Hartree
\end{lstlisting}

Orbitals and orbital energies are those of the KOBMP2 object; the singles
block and the couplings use the HF Fock matrix in that basis. Only the target
$k$-point is an external index, so no array with three $k$-point indices is
stored. With \code{symmetry=True} (default) the singles space is the
degenerate level of the band edge; the neglected coupling to the other bands
of the $k$-point changes the electron affinity of diamond by 0.01--0.03~eV.
The IP and EA agree with the EOM-KCCSD code of \pyscf{} (with $t_1=0$,
MP2-like $t_2$ and the same doubles--doubles block) to $10^{-9}$~Hartree.

\subsection{Comparing gaps}

Band gaps with $k$-point sampling depend strongly on the mesh and on the
treatment of the exchange divergence (\code{exxdiv} of the mean field). A
value on a single mesh is not a result. Series with \code{exxdiv=None} and
\code{exxdiv='ewald'} must be extrapolated, for example with
$a + bN_k^{-1/3} + cN_k^{-2/3}$, before they can be compared with experiment.
The script \code{benchmarks/diamond_gaps.py} computes the gaps of diamond for
HF and the variants above on one mesh.

\section{k-point downfolding: KOBDF}
\label{sec:kobdf}

\code{KOBDF(mp, ncas, nelecas, cas_list=None)}, or \code{mp.downfold(ncas,
nelecas)}, builds an active-space Hamiltonian from a \code{KOBMP2} object,
per unit cell,
\begin{equation}
  \hat H_{\mathrm{act}} = E_{\mathrm{core}}
  + \frac1{\Nk}\sum_{k}\sum_{tu}\bar h_{tu}(k)\,\hat a^\dagger_{tk}\hat a_{uk}
  + \frac1{2\Nk}\sum_{k_1k_2k_3}\sum_{tuvw}
  h^{(2)}_{tuvw}(k_1,k_2,k_3)\,
  \hat a^\dagger_{tk_1}\hat a^\dagger_{vk_3}\hat a_{wk_4}\hat a_{uk_2},
\end{equation}
with $\bar h(k)=h(k)+v^{\mathrm{HF}}_{\mathrm{core}}(k)+v[T](k)-v[T_{\mathrm{act}}](k)$
and $h^{(2)}=(tk_1,uk_2|vk_3,wk_4)/\Nk$. The cost is one KOBMP2 cycle plus
one cycle restricted to the active bands. With
\code{KOBMP2(mf).init_from_scf().downfold(...)} the downfolding is done at
second order on HF orbitals.

\begin{lstlisting}[style=python]
mp = KOBMP2(mf).run()
ham = mp.downfold(ncas=8, nelecas=8)
h1, e_core = ham.get_h1eff()          # (nkpts, ncas, ncas), constant per cell
ham.get_cderi(k1, k2)                 # three-index integrals of the active bands
ham.get_h2eff(k1, k2, k3)             # (t k1, u k2 | v k3, w k4) / Nk
ham.scf(); ham.mo_energy; ham.band_gap()      # mean field of the model
ham.self_energy(nocc=1, nvir=1)
ham.quasiparticle_gap('dynamic')      # second-order gap of the model
h1s, h2s, e0, nelec = ham.supercell() # real Hamiltonian of the supercell
ham.as_scf()                          # PySCF RHF object (CCSD, CASCI, DMRG)
ham.kernel()                          # FCI of the supercell, energy per cell
\end{lstlisting}

\begin{description}[style=nextline,leftmargin=3.2cm,labelwidth=3cm]
  \item[\code{ncas}] number of active bands per $k$-point.
  \item[\code{nelecas}] number of active electrons per unit cell (even).
  \item[\code{cas_list}] 0-based band indices, the same at every $k$-point.
        By default the \code{ncas} bands around the Fermi level.
  \item[\code{conv_tol}, \code{max_cycle}] convergence of \code{scf()}
        (defaults $10^{-10}$ and 50).
  \item[\code{fcisolver}] solver of \code{kernel()}.
\end{description}

\paragraph{Mean field and quasiparticles of the model.} \code{scf()} solves
Hartree--Fock for $\hat H_{\mathrm{act}}$; its bands contain the correlation
with the orbitals outside the active space statically, through
$v^{\mathrm{ext}}$. \code{self_energy()} adds the correlation within the
active space with its frequency dependence, and
\code{quasiparticle(scheme)} solves for the bands as in \code{kqp}.

\paragraph{Many-body solvers.} \code{supercell()} returns real integrals of
the Born--von K\'arm\'an supercell of $\Nk$ cells ($\Nk n_{\mathrm{cas}}$
orbitals). The real orbitals are obtained from time reversal and are not
localized. \code{as_scf()} wraps the integrals in a \pyscf{} \code{RHF}
object; its energies refer to the supercell. \code{kernel()} runs FCI and is
only feasible for small models.

\paragraph{Sum rule.} \code{energy_tot()} over the reference determinant
equals $E_{\mathrm{KOBMP2}}-E_c[T_{\mathrm{act}}]$; the second term is stored
as \code{e_corr_act}.

\paragraph{Madelung term.} With \code{exxdiv='ewald'} the constant contains
$-\xi N_{\mathrm{act}}/2$ for the active electrons, and \code{get_fock}
shifts the occupied levels by $-\xi$, as \pyscf{} does for KRHF; the
integrals do not contain the term.

From \code{examples/experimental/kobdf.py} (diamond, gth-szv,
$2\times2\times1$, \code{exxdiv=None}, six active bands):

\begin{lstlisting}[style=output]
E(KRHF) = -9.26630050   E(KOBMP2) = -9.41369394   gap of the KOBMP2 bands = 0.37745
h1eff: (4, 6, 6)   E_core = -12.34762518
sum rule: <H_act> - [E(KOBMP2) - E_c(T_act)] = 8.2e-13
E(HF of H_act) = -9.34556626   gap = 0.32196
second-order gap of H_act, static  = 0.30080
second-order gap of H_act, linear  = 0.30205
second-order gap of H_act, dynamic = 0.30205
supercell: 24 orbitals, 24 electrons; E(HF) = -9.34556626, E(CCSD) = -9.39969464 per cell
\end{lstlisting}

\section{Unrestricted k-point OBMP2: KUOBMP2}

\code{KUOBMP2} is the four-index code of version 0.2, moved into the package
largely unchanged; its interface is that of version 0.2 (\code{thresh},
\code{niter}, \code{shift}). For a closed-shell cell it does not reproduce
KOBMP2 (see Section~\ref{sec:known-issues}), it has the cost problems of the
old code, and its ionization analysis uses hard-coded band indices. It should
be used only for comparison with old results.
